\chapter{等离子体基础诊断}
\label{chap:diagnostics}

\begin{abstract}
等离子体是"看不见摸得着"的物质第四态。它不像固体、液体或普通气体那样可以直接观察，却又遍布从实验室到宇宙空间的各个角落。本章将系统介绍四种基础诊断手段——电探针、微波、光学光谱和磁探针——并阐明一个核心思想：\emph{诊断是等离子体理论预言与实验观测之间的桥梁}。每一种诊断方法，本质上都是将第2章的德拜屏蔽、第5章的磁流体力学、第6章的电磁波传播和第7章的动理学分布函数等抽象概念，转化为电压表、示波器和光谱仪上可读取的信号。
\end{abstract}

\section{引言：为什么需要诊断？}
\label{sec:12.1}

如果你走进一间等离子体实验室，首先注意到的可能不是等离子体本身，而是围绕真空腔体布置的纷繁复杂的探针、线圈、激光器和光谱仪。原因很简单：\emph{绝大多数等离子体并不直接可见}。低温等离子体可能发出微弱的辉光，但高温等离子体（如聚变装置中的等离子体）辐射的主要波段在不可见的X射线或紫外区；磁化等离子体中的物理过程更被束缚在复杂的磁场拓扑中，肉眼完全无法分辨。

\begin{definition}{等离子体诊断}{}{}
等离子体诊断（Plasma Diagnostics）是指利用电、磁、电磁波、粒子或光学手段，对等离子体的密度、温度、电位、成分、电场、磁场等宏观与微观参数进行\emph{非破坏性或可控扰动性}测量的实验技术体系。
\end{definition}

\begin{insight}{间接测量的本质}{}{}
诊断的核心逻辑可以概括为：\textbf{理论预言} $\rightarrow$ \textbf{物理效应} $\rightarrow$ \textbf{可测量信号} $\rightarrow$ \textbf{反演参数}。

例如，理论告诉我们电子存在特征振荡频率 $\omega_{pe}$（见第6章式\eqref{eq:plasma-frequency}）。当一束微波穿过等离子体时，如果微波频率 $\omega$ 接近 $\omega_{pe}$，传播特性就会发生剧烈变化——这个\emph{可测量}的相位移动或幅度衰减，就让我们得以反演出等离子体密度。诊断物理学家从不"直接看到"电子，而是通过理解电子的\emph{集体响应}来推断其存在和状态。
\end{insight}

按照测量原理，等离子体诊断可分为四大类：
\begin{enumerate}
    \item \textbf{电探针诊断}：将金属电极插入等离子体，通过收集电流-电压特性来测量局部参数；
    \item \textbf{电磁波诊断}：利用微波或激光在等离子体中的传播、散射或辐射特性；
    \item \textbf{光学光谱诊断}：分析等离子体自发辐射或散射光的谱线特征；
    \item \textbf{粒子与磁探针诊断}：测量逃逸粒子或感应磁场，反推等离子体电流与压强。
\end{enumerate}

本章将介绍这四种方法中最基础、最经典的原理，为后续学习更高级的诊断技术（如汤姆逊散射成像、电荷交换复合光谱、偏振干涉仪等）奠定物理图像基础。

\begin{tip}
\textbf{学习提示：}
学习诊断学时，建议始终带着两个问题阅读：\textbf{(1)} 这种诊断手段依赖的是等离子体的哪一个特征尺度或特征频率？（如德拜长度 $\lambda_D$、等离子体频率 $\omega_{pe}$、回旋频率 $\omega_{ce}$）\textbf{(2)} 该测量是\emph{局域}的（点测量）还是\emph{积分}的（线测量）？这两个问题将帮助你判断不同诊断方法的适用场景和局限性。
\end{tip}


\section{朗缪尔探针（Langmuir Probe）}
\label{sec:12.2}

朗缪尔探针是最古老、最简洁、也是信息密度最高的等离子体诊断工具之一。它只需一根细金属丝，插入等离子体中，通过扫描探针电压并测量收集电流，就能得到局部的电子温度 $T_e$、电子密度 $n_e$ 和等离子体电位 $V_p$。

\subsection{探针结构与基本电路}

根据几何构型和测量目的，朗缪尔探针分为三种基本类型：

\begin{definition}{单探针、双探针与发射探针}{}{}
\begin{itemize}
    \item \textbf{单探针（Single Probe）}：一根金属电极相对于\emph{真空腔壁}（或参考电极）扫描电压，测量电流。适用于电位已知的接地或壁面参考系统。
    \item \textbf{双探针（Double Probe）}：两根相同探针相互悬浮，在两探针之间施加扫描电压。不依赖腔壁电位，特别适合射频放电或绝缘腔体中的测量。
    \item \textbf{发射探针（Emissive Probe）}：通过加热使探针发射热电子，其电流特性会显著改变悬浮电位附近的曲线形态，可更精确地测量等离子体电位 $V_p$。
\end{itemize}
\end{definition}

单探针的等效电路极为简单：直流电源通过扫描电路向探针施加电压 $V$，串联电流表测量探针收集电流 $I$。探针通常选用耐高温、低溅射的金属（如钨、钼或铂），直径 $0.1\sim 1\,\mathrm{mm}$，长度 $1\sim 10\,\mathrm{mm}$。

\subsection{I-V特性曲线与四个特征区域}

朗缪尔探针的电流-电压（$I$-$V$）曲线是诊断的\emph{全部信息来源}。其典型形状如图\ref{fig:langmuir-iv}所示，可分为四个特征区域。

\begin{theorem}{朗缪尔探针收集电流}{}{}
对于无磁化的、麦克斯韦分布的等离子体，单探针的收集电流满足
\begin{equation}
I(V) = I_{\mathrm{i,sat}} - I_{\mathrm{e,sat}}\exp\!\left(\frac{V - V_p}{T_e}\right),
\label{eq:langmuir-iv}
\end{equation}
其中 $V_p$ 为等离子体电位，$T_e$ 以能量单位（$\mathrm{eV}$）表示，
\begin{align}
I_{\mathrm{i,sat}} &\approx \frac{1}{2} e n_i A \bar{v}_i \quad \text{（离子饱和电流）},
\label{eq:ion-sat}\\
I_{\mathrm{e,sat}} &\approx \frac{1}{4} e n_e A \bar{v}_e \quad \text{（电子饱和电流）},
\label{eq:elec-sat}
\end{align}
$A$ 为探针收集面积，$\bar{v}_\alpha = \sqrt{8 T_\alpha / (\pi m_\alpha)}$ 为平均热速度。
\end{theorem}

\begin{insight}{为什么I-V曲线呈"S"形？}{}{}
探针本质上是一个可调控电位的\emph{带电粒子收集器}。理解曲线的关键是：电子和离子带相反电荷，但电子质量远小于离子，因此电子热速度 $\bar{v}_e$ 远大于离子热速度 $\bar{v}_i$。

\begin{itemize}
    \item \textbf{远负电位区}：探针电位远低于等离子体电位，电子被排斥，只有离子被收集。此时电流饱和为离子饱和电流 $I_{\mathrm{i,sat}}$（小负值）。
    \item \textbf{过渡区}：电位逐渐升高，越来越多的电子克服势垒到达探针。由于电子速度服从玻尔兹曼分布，电流随电压呈指数增长。
    \item \textbf{远正电位区}：探针电位高于等离子体电位，离子被排斥，电子被大量收集。但由于探针周围形成电子耗尽鞘层，电流不会无限增长，而是趋于电子饱和电流 $I_{\mathrm{e,sat}}$（大正值）。
\end{itemize}

曲线的"S"形，本质上就是等离子体中正负电荷响应\emph{不对称性}的直观体现。
\end{insight}

\begin{warning}{探针尺寸与德拜长度的关系}{}{}
\textbf{误区：}探针尺寸必须远大于德拜长度 $\lambda_D$（见第2章式\eqref{eq:debye-length}）才能工作。

\textbf{正解：}恰恰相反。如果探针尺寸远大于 $\lambda_D$，探针会严重扰动等离子体，形成厚的非中性鞘层，理论分析将极为复杂。理想情况下，探针尺寸应\emph{与} $\lambda_D$ 相当或更小，这样鞘层较薄，空间电荷效应可控，式\eqref{eq:langmuir-iv}的薄鞘理论才能适用。但探针也不能太小，否则收集电流过弱，信噪比不足。实际实验中，探针半径 $r_p \sim (1\sim 10)\,\lambda_D$ 是常见折中。
\end{warning}

\subsection{电子温度 $T_e$ 的求法：过渡区斜率法}

在过渡区，离子电流可忽略，$I \approx I_{\mathrm{e,sat}}\exp[(V - V_p)/T_e]$。取对数后：
\begin{equation}
\ln I = \ln I_{\mathrm{e,sat}} + \frac{V - V_p}{T_e}.
\end{equation}

\begin{corollary}{电子温度的斜率公式}{}{}
在过渡区，对 $\ln I$ 关于 $V$ 求导可得
\begin{equation}
\frac{\diff (\ln I)}{\diff V} = \frac{1}{T_e}
\quad\Longrightarrow\quad
T_e = \left[\frac{\diff (\ln I)}{\diff V}\right]^{-1}.
\label{eq:te-slope}
\end{equation}
因此，$\ln I$-$V$ 曲线的斜率倒数即给出电子温度 $T_e$（单位为 $\mathrm{eV}$）。
\end{corollary}

\begin{insight}{斜率为何反比于温度？}{}{}
式\eqref{eq:te-slope}的物理图像非常直观。电子温度越高，热运动越剧烈，能够克服探针负电位"势垒"到达探针的电子就越多。这意味着温度越高，收集电流对电位变化越"不敏感"——你需要改变更大的电位才能显著改变到达探针的电子数。因此，$T_e$ 越大，$\ln I$-$V$ 曲线越平缓，斜率越小，其倒数（温度）就越大。反之，低温等离子体中曲线非常陡峭。
\end{insight}

\subsection{电子密度 $n_e$ 的求法}

一旦从斜率法求得 $T_e$，就可以利用饱和区电流反推密度。

\begin{corollary}{电子密度公式}{}{}
利用电子饱和电流式\eqref{eq:elec-sat}，电子密度为
\begin{equation}
n_e = \frac{4 I_{\mathrm{e,sat}}}{e A \bar{v}_e}
= \frac{4 I_{\mathrm{e,sat}}}{e A}\sqrt{\frac{\pi m_e}{8 T_e}}
= I_{\mathrm{e,sat}} \sqrt{\frac{2\pi m_e}{e^2 A^2 T_e}}.
\label{eq:ne-from-sat}
\end{equation}

对于低温等离子体（$T_e \lesssim T_i$），电子饱和电流区难以精确确定，更常用\textbf{离子饱和电流}估算：
\begin{equation}
n_e \approx n_i \approx \frac{2 I_{\mathrm{i,sat}}}{e A \bar{v}_i}
= \frac{2 I_{\mathrm{i,sat}}}{e A}\sqrt{\frac{\pi m_i}{8 T_i}}.
\label{eq:ne-from-ion-sat}
\end{equation}
\end{corollary}

\begin{warning}{磁化等离子体中的探针}{}{}
在存在磁场的等离子体（如磁约束聚变装置）中，电子回旋半径 $r_{ce} = v_{\perp e}/\omega_{ce}$（见第5章）可能远小于探针尺寸。此时电子沿磁场自由运动，但在垂直方向被磁化，探针收集电流不再由式\eqref{eq:langmuir-iv}简单描述。电子饱和电流被磁化抑制，过渡区斜率也不再单纯给出 $T_e$。磁化探针分析需要引入\emph{磁化鞘层理论}，通常需要数值模拟或经验标定，不能直接套用无磁化公式。
\end{warning}

\subsection{电子能量分布函数（EEDF）}

如果等离子体中的电子分布偏离麦克斯韦分布，仅靠 $T_e$ 一个参数无法完整描述。朗缪尔探针的一个重要扩展功能是通过 $I$-$V$ 曲线求出完整的电子能量分布函数（EEDF）。

\begin{theorem}{Druyvesteyn公式}{}{}
对于各向同性的电子速度分布，探针收集电流的二阶导数与EEDF $f(\eps)$ 的关系为
\begin{equation}
\frac{\diff^2 I}{\diff V^2} = \frac{e^3 A}{2 m_e} f(-eV)
\quad \text{（当 } V < V_p \text{ 时）}.
\label{eq:druyvesteyn}
\end{equation}
其中 $\eps = -e(V - V_p)$ 为电子相对于等离子体电位的能量，$f(\eps)$ 为按能量归一化的分布函数。
\end{theorem}

\begin{insight}{为什么二阶导数给出分布函数？}{}{}
探针收集电流是\emph{所有能量高于势垒}的电子的通量积分：$I \propto \int_{-eV}^\infty f(\eps) v(\eps) \,\diff\eps$。对 $V$ 求一阶导数消去了积分下限的被积函数（莱布尼茨法则），再求二阶导数就提取出了分布函数本身。这相当于在说：\emph{电流曲线在某处的"弯曲程度"，正比于具有该能量的电子数量}。这一结果最早由Druyvesteyn在1930年严格导出，是探针诊断从"测温度"到"测分布"的理论飞跃。
\end{insight}

\begin{example}{数值估算：实验室氩等离子体参数}{}{}
\difficulty{基础}
某直流辉光放电氩等离子体中，朗缪尔探针（圆柱形，半径 $r = 0.25\,\mathrm{mm}$，长度 $L = 5\,\mathrm{mm}$）测得离子饱和电流 $I_{\mathrm{i,sat}} = 0.15\,\mathrm{mA}$，$\ln I$-$V$ 曲线过渡区斜率为 $3.5\,\mathrm{V}^{-1}$。估算 $T_e$ 和 $n_e$。

\textbf{解：}
\begin{enumerate}
    \item 电子温度：由式\eqref{eq:te-slope}，$T_e = 1/3.5 \approx 0.29\,\mathrm{V} = 0.29\,\mathrm{eV}$（约 $3300\,\mathrm{K}$）。
    \item 探针面积：$A = 2\pi r L + 2\pi r^2 \approx 2\pi \times 0.25\times 10^{-3}\times 5\times 10^{-3} \approx 7.85\times 10^{-6}\,\mathrm{m}^2$（忽略端面）。
    \item 离子密度：对于氩离子 $m_i = 40\,\mathrm{amu} \approx 6.64\times 10^{-26}\,\mathrm{kg}$，假设 $T_i \approx T_e = 0.29\,\mathrm{eV}$，则 $\bar{v}_i = \sqrt{8T_i/(\pi m_i)} \approx 2.1\times 10^3\,\mathrm{m/s}$。由式\eqref{eq:ne-from-ion-sat}：
    \begin{equation}
    n_e \approx \frac{2\times 0.15\times 10^{-3}}{1.6\times 10^{-19}\times 7.85\times 10^{-6}\times 2.1\times 10^3} \approx 1.1\times 10^{17}\,\mathrm{m}^{-3}.
    \end{equation}
\end{enumerate}
这是一个典型的低气压辉光放电密度，德拜长度 $\lambda_D \approx 740\,\mu\mathrm{m}$，探针半径 $r = 250\,\mu\mathrm{m} \approx 0.34\,\lambda_D$，鞘层较薄，薄鞘理论尚可适用。
\end{example}

\subsection{悬浮电位与历史回顾}

\begin{definition}{悬浮电位}{}{}
当探针与外部电路断开（净电流 $I = 0$）时，探针自发达到的电位称为\textbf{悬浮电位}（Floating Potential），记作 $V_f$。由 $I_{\mathrm{i,sat}} = I_{\mathrm{e,sat}}\exp[(V_f - V_p)/T_e]$，可得
\begin{equation}
V_f - V_p = T_e \ln\!\left(\frac{I_{\mathrm{i,sat}}}{I_{\mathrm{e,sat}}}\right)
= T_e \ln\!\left(\sqrt{\frac{m_e}{m_i}}\right)
\approx -T_e \ln\!\left(\sqrt{\frac{m_i}{2.3\,\mathrm{MeV}/c^2}}\right).
\label{eq:floating-potential}
\end{equation}
\end{definition}

\begin{insight}{为什么悬浮电位是负的？}{}{}
式\eqref{eq:floating-potential}中，$V_f < V_p$ 的根本原因在于\textbf{电子的流动性远高于离子}。电子质量小，热速度大，"冲向"探针的通量天然就比离子大几个数量级。如果探针保持中性电位，它会被\emph{过度充电}为负电位。只有探针自身积累足够负的电荷，形成排斥电子的势垒，才能使电子通量降低到与离子通量相等，达到动态平衡。对于氢等离子体，$V_f - V_p \approx -3.0\,T_e$；对于氩等离子体，这一偏差可达 $-4.7\,T_e$。因此，悬浮电位\emph{不是}等离子体电位，这是初学者最容易混淆的概念。
\end{insight}

\begin{history}{朗缪尔与探针诊断的诞生}{}{}
1924年，美国通用电气研究实验室的欧文·朗缪尔（Irving Langmuir，1881--1957）在研究汞蒸气放电时，首次系统地利用插入金属丝测量电流-电压特性，从而推导出电子温度和密度。朗缪尔当时并没有复杂的等离子体理论框架，他凭借对气体放电的深刻直觉，提出了鞘层（Sheath）的概念，并正确认识到探针周围必然存在一个非中性的空间电荷区。这项工作不仅为等离子体诊断奠定了实验基础，也直接催生了"等离子体"（Plasma）一词在物理学中的现代用法——朗缪尔本人于1928年正式提议用这个词描述电离气体。朗缪尔探针至今仍是等离子体实验室的"标配"，其简洁之美历经百年而不衰。
\end{history}


\section{微波干涉诊断}
\label{sec:12.3}

当等离子体密度足够高时，它会对电磁波的传播产生显著影响。微波干涉诊断正是利用这一效应，通过测量微波穿过等离子体后的相位变化，来反演出等离子体柱的平均电子密度。

\subsection{基本原理：等离子体作为色散介质}

回顾第6章，电磁波在无碰撞冷等离子体中的色散关系为
\begin{equation}
\omega^2 = \omega_{pe}^2 + k^2 c^2
\quad\Longrightarrow\quad
n_r = \frac{kc}{\omega} = \sqrt{1 - \frac{\omega_{pe}^2}{\omega^2}},
\label{eq:dispersion-cold}
\end{equation}
其中 $\omega_{pe} = \sqrt{n_e e^2 / (\eps_0 m_e)}$ 为等离子体频率，$n_r$ 为折射率。

\begin{definition}{截止密度}{}{}
当探测频率 $\omega$ 等于等离子体频率时，折射率 $n_r = 0$，波无法传播。由此定义\textbf{截止密度}（Cutoff Density）：
\begin{equation}
n_{\mathrm{cut}} = \frac{\eps_0 m_e \omega^2}{e^2}
\approx 1.24\times 10^{14}\,\left(\frac{f}{\mathrm{GHz}}\right)^2\,\mathrm{m}^{-3}.
\label{eq:cutoff-density}
\end{equation}
当 $n_e > n_{\mathrm{cut}}$ 时，微波将被反射而无法穿透等离子体。
\end{definition}

\begin{insight}{折射率为什么小于1？}{}{}
等离子体中，自由电子在外加电场作用下做受迫振荡，其振荡相位与驱动电场相差 $\pi$（电子带负电，加速度与电场反向）。电子振荡产生的感应电场与外加电场相消，导致有效介电常数 $\eps_{\mathrm{eff}} = \eps_0(1 - \omega_{pe}^2/\omega^2)$ 小于真空介电常数 $\eps_0$。因此等离子体的折射率 $n_r = \sqrt{\eps_{\mathrm{eff}}/\eps_0} < 1$。波在等离子体中的相速度 $v_p = c/n_r > c$，但群速度 $v_g = c n_r < c$，能量传播变慢。这正是波穿过等离子体后产生相位延迟的物理根源。
\end{insight}

\subsection{相位移动与密度关系}

在实际诊断中，通常选择 $\omega \gg \omega_{pe}$（即 $n_r \approx 1$），此时微波可以穿透等离子体，但会产生微小的相位移动。对式\eqref{eq:dispersion-cold}作展开：
\begin{equation}
n_r \approx 1 - \frac{1}{2}\frac{\omega_{pe}^2}{\omega^2}
= 1 - \frac{e^2}{2\eps_0 m_e \omega^2} n_e.
\end{equation}

\begin{theorem}{微波干涉相位公式}{}{}
微波沿路径 $L$ 穿过等离子体后，相对于真空路径的相位移动为
\begin{equation}
\Delta\phi = \frac{\omega}{c}\int_0^L (1 - n_r)\,\diff l
= \frac{e^2}{2\eps_0 m_e c \omega} \int_0^L n_e\,\diff l.
\label{eq:phase-shift}
\end{equation}

对于均匀等离子体柱（密度 $n_e$，长度 $L$），上式简化为
\begin{equation}
\Delta\phi = \frac{e^2}{2\eps_0 m_e c \omega}\, n_e L
= 8.97\times 10^{-17}\,\frac{\lambda_0}{\mathrm{m}} \,\frac{n_e}{\mathrm{m}^{-3}}\,\frac{L}{\mathrm{m}}\,\mathrm{rad},
\label{eq:phase-shift-uniform}
\end{equation}
其中 $\lambda_0 = 2\pi c/\omega$ 为真空波长。
\end{theorem}

\begin{insight}{为什么相位移动正比于密度和长度？}{}{}
式\eqref{eq:phase-shift}具有极其简洁的物理意义：每一个电子对电磁波相位延迟的"贡献"是独立的，且贡献量与电子密度成正比。等离子体柱越长，参与相互作用的电子越多，累积相位移动越大。值得注意的是，\textbf{离子基本不参与}，因为离子质量大，在微波频率下几乎不动，对介电常数的修正被 $(m_e/m_i)$ 因子抑制。微波干涉诊断本质上是一种\textbf{纯电子密度测量}。
\end{insight}

\begin{warning}{相位移动与频率的关系}{}{}
\textbf{误区：}使用更高频率的微波可以得到更大的相位移动，从而提高灵敏度。

\textbf{正解：}由式\eqref{eq:phase-shift}，$\Delta\phi \propto 1/\omega$。频率越高，相位移动反而\emph{越小}。但高频率允许测量更高密度（截止密度 $n_{\mathrm{cut}} \propto f^2$）。因此，实验设计中需要权衡：频率太低则无法穿透高密度等离子体；频率太高则相位移动微弱，难以检测。典型选择是使 $n_{\mathrm{cut}}$ 为预期最大密度的 $2\sim 3$ 倍。例如，对于峰值密度 $n_e \sim 10^{18}\,\mathrm{m}^{-3}$，选用 $f = 50\,\mathrm{GHz}$（$n_{\mathrm{cut}} \approx 3.1\times 10^{18}\,\mathrm{m}^{-3}$）是合理的。
\end{warning}

\begin{example}{数值估算：托卡马克芯部密度测量}{}{}
\difficulty{基础}
某托卡马克装置使用波长 $\lambda_0 = 2\,\mathrm{mm}$（$f = 150\,\mathrm{GHz}$）的微波进行干涉测量，等离子体柱直径（微波穿越路径）约 $L = 1.0\,\mathrm{m}$。若测得相位移动 $\Delta\phi = 10\pi\,\mathrm{rad}$（5个条纹移动），估算平均电子密度。

\textbf{解：}由式\eqref{eq:phase-shift-uniform}，
\begin{equation}
n_e = \frac{2\eps_0 m_e c \omega \Delta\phi}{e^2 L}
= \frac{2\eps_0 m_e c^2 \Delta\phi}{e^2 L} \cdot \frac{2\pi}{\lambda_0}.
\end{equation}
代入数值：$\eps_0 = 8.85\times 10^{-12}\,\mathrm{F/m}$，$m_e = 9.11\times 10^{-31}\,\mathrm{kg}$，$c = 3\times 10^8\,\mathrm{m/s}$，$e = 1.6\times 10^{-19}\,\mathrm{C}$，$\Delta\phi = 10\pi$，$\lambda_0 = 2\times 10^{-3}\,\mathrm{m}$，$L = 1.0\,\mathrm{m}$：
\begin{equation}
n_e \approx \frac{2\times 8.85\times 10^{-12}\times 9.11\times 10^{-31}\times (3\times 10^8)^2 \times 10\pi}{(1.6\times 10^{-19})^2 \times 1.0 \times 2\times 10^{-3}}
\approx 8.9\times 10^{18}\,\mathrm{m}^{-3}.
\end{equation}
这一密度处于微波干涉诊断的典型范围内（$10^{18}\sim 10^{20}\,\mathrm{m}^{-3}$）。对于更高密度（如聚变堆级 $10^{20}\,\mathrm{m}^{-3}$），需要使用亚毫米波或远红外激光干涉仪。
\end{example}


\section{激光汤姆逊散射（Laser Thomson Scattering）}
\label{sec:12.4}

如果说朗缪尔探针是"电接触式"的局部测量，微波干涉是"电磁波积分式"的线测量，那么激光汤姆逊散射（Laser Thomson Scattering, LTS）则是一种\emph{非接触、局域、可分辨}的精密诊断。它被誉为等离子体诊断的"黄金标准"，尤其适用于高温聚变等离子体。

\subsection{基本物理图像}

当一束强激光（波长通常为 $\mathrm{nm}$ 到 $\mu\mathrm{m}$ 量级）照射到等离子体上时，激光电场驱动自由电子做受迫振荡，电子再以偶极辐射的形式向各个方向散射电磁波。这就是汤姆逊散射。

\begin{definition}{汤姆逊散射微分截面}{}{}
单个自由电子对非偏振入射光的汤姆逊散射微分截面为
\begin{equation}
\frac{\diff\sigma}{\diff\Omega} = r_e^2 \sin^2\theta_p,
\label{eq:thomson-cross-section}
\end{equation}
其中 $r_e = e^2/(4\pi\eps_0 m_e c^2) \approx 2.82\times 10^{-15}\,\mathrm{m}$ 为经典电子半径，$\theta_p$ 为散射角（散射光与入射光电场方向的夹角）。对非偏振光在所有偏振方向平均后，得到常用的
\begin{equation}
\frac{\diff\sigma}{\diff\Omega} = \frac{r_e^2}{2}(1 + \cos^2\theta),
\label{eq:thomson-unpolarized}
\end{equation}
$\theta$ 为入射方向与散射方向的夹角。总截面为 $\sigma_T = 8\pi r_e^2/3 \approx 6.65\times 10^{-29}\,\mathrm{m}^2$。
\end{definition}

\subsection{非相干散射与电子温度测量}

汤姆逊散射的核心诊断价值在于：\textbf{散射光谱的宽度直接反映电子的热运动速度分布}。

\begin{theorem}{非相干汤姆逊散射光谱}{}{}
当电子密度满足 $n_e \ll n_{\mathrm{co}} \sim 10^{23}\,\mathrm{m}^{-3}$（非相干散射条件，见下文）时，不同电子的散射波相位随机，总散射强度为各电子贡献的非相干叠加。对于麦克斯韦分布的电子，散射光谱呈高斯形：
\begin{equation}
S(\omega_s) \propto \exp\!\left[-\frac{(\omega_s - \omega_i)^2}{2\Delta\omega^2}\right],
\label{eq:thomson-gaussian}
\end{equation}
其中 $\omega_i$ 为入射激光频率，光谱宽度（半高全宽，$1/\ee$ 宽度为 $2\Delta\omega$）满足
\begin{equation}
\frac{\Delta\omega}{\omega_i} = 2\sqrt{2\ln 2}\,\sin\frac{\theta}{2}\,\frac{v_{th}}{c}
= 2\sqrt{2\ln 2}\,\sin\frac{\theta}{2}\,\sqrt{\frac{2T_e}{m_e c^2}}.
\label{eq:thomson-width}
\end{equation}
\end{theorem}

\begin{insight}{为什么散射光谱的宽度就是温度？}{}{}
这是汤姆逊散射最优美的物理图像。电子以热速度 $v$ 随机运动，对入射激光产生多普勒频移：频率偏移 $\Delta\omega = \vect{k}_s\cdot\vect{v} - \vect{k}_i\cdot\vect{v} = \vect{q}\cdot\vect{v}$，其中 $\vect{q} = \vect{k}_s - \vect{k}_i$ 为散射波矢。由于 $|\vect{q}| = 2k_i\sin(\theta/2)$，而电子速度服从麦克斯韦分布（见第7章），速度投影 $q\cdot v$ 的分布也是高斯型，其宽度正比于热速度 $v_{th} = \sqrt{2T_e/m_e}$。因此，\textbf{测量散射光谱就是测量电子速度分布}，光谱宽度就是温度，光谱形状偏离高斯就是分布偏离麦克斯韦。这直接将动理学理论（第7章）的抽象分布函数 $f(\vect{v})$，转化为光谱仪上一条可测量的谱线。
\end{insight}

\subsection{为什么不测量离子？}

离子同样会受到激光电场作用，也会散射光。但汤姆逊散射诊断中几乎只观测电子散射信号，原因在于：

\begin{corollary}{离子与电子散射截面之比}{}{}
汤姆逊散射截面与粒子质量平方成反比：
\begin{equation}
\frac{\sigma_i}{\sigma_e} = \left(\frac{m_e}{m_i}\right)^2.
\label{eq:ion-electron-ratio}
\end{equation}
对于氢等离子体，$m_i/m_e \approx 1836$，因此离子散射截面仅为电子的 $3\times 10^{-7}$。对于氘等离子体，这一比例更小。此外，离子热速度低，多普勒展宽极小，其散射信号在光谱上浓缩为极窄的"离子成分峰"（通常位于中心频率附近），极易与电子宽光谱成分区分。
\end{corollary}

\begin{insight}{质量优势的极端性}{}{}
式\eqref{eq:ion-electron-ratio}揭示了自然界中一个深刻的数量级规律：由于电子质量远小于离子，电子在相同电场下的加速度大 $(m_i/m_e)$ 倍，辐射振幅大 $(m_i/m_e)$ 倍，辐射功率（正比于振幅平方）大 $(m_i/m_e)^2$ 倍。这导致汤姆逊散射几乎完全是电子的"独角戏"，离子只是背景噪音。这也是为什么汤姆逊散射是\emph{最纯净的电子温度诊断}——它几乎不受离子成分、杂质或电荷态的影响。
\end{insight}

\begin{warning}{相干散射与非相干散射的界限}{}{}
当密度升高时，德拜球内的电子数 $N_D = n_e \lambda_D^3$ 不再是大量，电子间的相关性增强，散射波不再完全随机叠加，而会出现集体效应——\textbf{相干散射}（或称为集体汤姆逊散射）。此时散射光谱不再单纯反映电子温度，还会出现与离子声波相关的窄峰（离子成分）和与电子等离子体波相关的宽峰。非相干散射的判据通常取 $k\lambda_D \gg 1$（$k = |\vect{q}|$），即探测尺度远小于德拜长度。对于 $90^\circ$ 散射和 $1064\,\mathrm{nm}$（Nd:YAG）激光，这一条件通常在 $n_e \ll 10^{23}\,\mathrm{m}^{-3}$ 时满足。聚变堆芯密度（$\sim 10^{20}\,\mathrm{m}^{-3}$）远低于此阈值，因此非相干近似仍然适用。
\end{warning}


\section{光谱诊断基础}
\label{sec:12.5}

等离子体是丰富的光源：从射频放电的淡紫色辉光，到弧光放电的耀眼白光，再到聚变等离子体的X射线辐射。光谱诊断通过分析这些光的颜色（波长）和强度，获取等离子体的温度、密度、成分和电场信息。

\subsection{线光谱：激发态温度与离子成分}

等离子体中的原子或离子被电子碰撞激发到高能级，退激发时辐射特征波长的光子，形成线光谱（Line Spectrum）。

\begin{definition}{激发温度与玻尔兹曼图}{}{}
在局部热力学平衡（LTE）假设下，激发态粒子数服从玻尔兹曼分布：
\begin{equation}
\frac{n_j}{n_i} = \frac{g_j}{g_i}\exp\!\left(-\frac{E_j - E_i}{T_{\mathrm{exc}}}\right),
\label{eq:boltzmann-excitation}
\end{equation}
其中 $g_j$、$g_i$ 为统计权重，$E_j$、$E_i$ 为能级能量。观测两条谱线的强度比 $I_{ji}/I_{lk}$，可反推出\textbf{激发温度} $T_{\mathrm{exc}}$：
\begin{equation}
\frac{I_{ji}}{I_{lk}} = \frac{A_{ji} g_j \lambda_{lk}}{A_{lk} g_l \lambda_{ji}} \exp\!\left(-\frac{E_j - E_l}{T_{\mathrm{exc}}}\right),
\label{eq:line-ratio}
\end{equation}
$A_{ji}$、$A_{lk}$ 为爱因斯坦自发辐射系数。
\end{definition}

\begin{insight}{谱线强度比为什么是"温度计"？}{}{}
式\eqref{eq:line-ratio}的本质是\textbf{能量杠杆}。两条谱线来自不同的激发能级，能级差 $\Delta E$ 就是这把"尺子"。温度低时，粒子大多集中在基态和低激发态，高能级谱线很弱；温度升高时，高能级被更多布居，高能谱线相对于低能谱线增强。通过测量强度比，我们实际上是在问："等离子体中的热涨落，有没有足够能量把粒子抬升到那个高度？" 如果 $T_{\mathrm{exc}} \ll \Delta E$，那条谱线几乎熄灭；如果 $T_{\mathrm{exc}} \sim \Delta E$，两条谱线势均力敌。玻尔兹曼因子 $\exp(-\Delta E/T)$ 将温度信息编码在光强中。
\end{insight}

此外，不同元素和电离态有其独特的指纹谱线。通过辨认谱线波长，可以诊断等离子体的\textbf{化学成分}和\textbf{电离度}——这对杂质控制（聚变装置）和工艺监控（等离子体刻蚀）至关重要。

\subsection{连续光谱：韧致辐射与电子温度}

除了线光谱，等离子体还产生连续光谱（Continuum），主要来源是\textbf{韧致辐射}（Bremsstrahlung，德文"制动辐射"）：自由电子经过离子附近时被库仑场偏转，产生加速辐射。

\begin{theorem}{韧致辐射功率谱}{}{}
对于非相对论性麦克斯韦电子分布，单位体积、单位频率、单位立体角的韧致辐射功率谱为
\begin{equation}
\frac{\diff P}{\diff V\,\diff\omega\,\diff\Omega} = \frac{e^6}{12\pi^3\eps_0^3 c^3 m_e^2}\sqrt{\frac{2m_e}{3\pi T_e}}\,n_e n_i Z^2\,\exp\!\left(-\frac{\hbar\omega}{T_e}\right) \bar{g}(\omega, T_e),
\label{eq:bremsstrahlung}
\end{equation}
其中 $Z$ 为离子电荷数，$\bar{g}$ 为冈特因子（Gaunt factor，量级为1的修正因子）。在光学薄等离子体中，连续谱强度的指数衰减特征 $\propto \exp(-\hbar\omega/T_e)$ 直接给出电子温度。
\end{theorem}

\begin{insight}{韧致辐射是"温度计"还是"密度计"？}{}{}
式\eqref{eq:bremsstrahlung}表明，韧致辐射功率同时正比于 $n_e n_i$（密度信息）和 $T_e^{-1/2}\exp(-\hbar\omega/T_e)$（温度信息）。但在不同频率区间，诊断侧重不同：在低频区（$\hbar\omega \ll T_e$），指数因子接近1，辐射强度主要取决于密度；在高频区（$\hbar\omega \sim T_e$），指数衰减敏感地依赖于温度。因此，通过测量连续谱的"斜率"（对数强度对频率的导数），可以提取 $T_e$，而绝对强度则与密度有关。这再次体现了诊断的多参数耦合特性——单一的测量通常无法同时确定所有参数，需要结合多个诊断手段。
\end{insight}

\subsection{斯塔克展宽：谱线里的密度信息}

等离子体中的带电粒子产生微观电场，这些电场通过线性或二次斯塔克效应，使发射或吸收谱线发生展宽和位移。这就是\textbf{斯塔克展宽}（Stark Broadening）。

\begin{definition}{斯塔克展宽与密度关系}{}{}
对于氢原子及其类氢离子，线性斯塔克效应占主导，谱线展宽 $\Delta\lambda_{1/2}$（半高全宽）与电子密度的关系近似为
\begin{equation}
\Delta\lambda_{1/2} \approx C\,n_e^{2/3},
\label{eq:stark-broadening}
\end{equation}
其中 $C$ 为与谱线种类和等离子体温度有关的系数。对于氢的 $H_\alpha$ 线（$656.3\,\mathrm{nm}$），在 $n_e \sim 10^{20}\sim 10^{24}\,\mathrm{m}^{-3}$、$T_e \sim 1\sim 10\,\mathrm{eV}$ 范围内，典型的经验公式为
\begin{equation}
\Delta\lambda_{1/2}\,\big[\mathrm{nm}\big] \approx 2.5\times 10^{-17}\left(\frac{n_e}{\mathrm{m}^{-3}}\right)^{2/3}.
\end{equation}
\end{definition}

\begin{insight}{为什么展宽正比于 $n_e^{2/3}$？}{}{}
线性斯塔克效应中，能级位移 $\Delta E$ 正比于外加电场 $\mathcal{E}$。等离子体中，带电粒子产生的电场在距离 $r$ 处为 $\mathcal{E} \sim e/(4\pi\eps_0 r^2)$。而最近的带电粒子典型距离由密度决定：$r \sim n_e^{-1/3}$（平均间距）。因此典型电场 $\mathcal{E}_{\mathrm{typ}} \sim e n_e^{2/3}/(4\pi\eps_0)$，能级位移和由此导致的谱线展宽也就正比于 $n_e^{2/3}$。这个简单的标度关系是光谱法测量密度的理论基础。需要注意的是，斯塔克展宽对温度也有弱依赖，因为离子热运动产生的微场统计分布随温度变化。精确的密度反演通常需要查阅专业数据库（如STARK-B）。
\end{insight}

\begin{example}{数值估算：电弧等离子体的 $H_\alpha$ 展宽}{}{}
\difficulty{基础}
某氢电弧等离子体中，光谱仪测得 $H_\alpha$ 线（$656.3\,\mathrm{nm}$）的半高全宽为 $\Delta\lambda_{1/2} = 2.5\,\mathrm{nm}$。假设温度在 $1\sim 10\,\mathrm{eV}$ 范围内，粗略估算电子密度。

\textbf{解：}由近似标度关系 $n_e \approx (\Delta\lambda_{1/2} / 2.5\times 10^{-17})^{3/2}$，代入 $\Delta\lambda_{1/2} = 2.5\times 10^{-9}\,\mathrm{m}$：
\begin{equation}
n_e \approx \left(\frac{2.5\times 10^{-9}}{2.5\times 10^{-17}}\right)^{3/2}
= (10^8)^{3/2} = 10^{12}\,\mathrm{m}^{-3}.
\end{equation}
等等，这个数值明显偏小。实际上对于氢 $H_\alpha$，常用经验公式在 $n_e \sim 10^{20}\,\mathrm{m}^{-3}$ 时 $\Delta\lambda_{1/2} \sim 0.1\,\mathrm{nm}$ 量级；$2.5\,\mathrm{nm}$ 的展宽对应 $n_e \sim 10^{23}\,\mathrm{m}^{-3}$ 量级。这说明\textbf{数值估算时必须使用正确的系数}。这个例子恰恰说明：光谱诊断的定量反演严重依赖精确的原子物理数据（如STARK-B数据库），标度关系只能用于量级估计。实际上，若使用更精确的公式：
\begin{equation}
n_e \approx 10^{23}\,\mathrm{m}^{-3}\times\left(\frac{\Delta\lambda_{1/2}}{0.5\,\mathrm{nm}}\right)^{3/2} \approx 10^{23}\times(5)^{3/2} \approx 1.1\times 10^{24}\,\mathrm{m}^{-3},
\end{equation}
对应于高电流密度电弧的芯部密度。
\end{example}


\section{磁探针（Magnetic Probe / Rogowski Coil）}
\label{sec:12.6}

前几节讨论的诊断手段主要测量等离子体的微观参数（$n_e$、$T_e$、分布函数）。但等离子体作为导电流体，其宏观行为（电流、磁场、压强）同样重要——这些正是第5章磁流体力学（MHD）的核心变量。磁探针和Rogowski线圈就是测量等离子体宏观电磁状态的基础工具。

\subsection{磁探针原理：法拉第感应}

\begin{definition}{磁探针}{}{}
\textbf{磁探针}（Magnetic Probe，又称磁场拾取线圈）是将小型线圈插入或靠近等离子体，通过测量感应电动势来探测局部磁场 $\vect{B}$ 或其时间变化率的诊断装置。
\end{definition}

根据法拉第电磁感应定律（见第5章），
\begin{equation}
\oint \vect{E}\cdot\diff\vect{l} = -\frac{\diff\Phi_B}{\diff t},
\label{eq:faraday}
\end{equation}
其中磁通量 $\Phi_B = N \int \vect{B}\cdot\diff\vect{S}$，$N$ 为线圈匝数。线圈两端的开路电压为
\begin{equation}
\mathcal{E} = -N A \frac{\diff B}{\diff t},
\label{eq:magnetic-probe}
\end{equation}
$A$ 为单匝面积，假设磁场垂直穿过线圈平面。对电压进行时间积分即可得到磁场变化：
\begin{equation}
\Delta B = -\frac{1}{NA}\int \mathcal{E}\,\diff t.
\label{eq:b-integration}
\end{equation}

\begin{insight}{为什么是电压的积分给出磁场？}{}{}
式\eqref{eq:magnetic-probe}告诉我们，磁探针的\emph{直接输出}是电压，而电压正比于磁场的\emph{变化率}，不是磁场本身。这源于电磁感应的物理本质：变化的磁场产生涡旋电场，驱动线圈中的电荷运动。对于稳恒磁场（$\diff B/\diff t = 0$），理想线圈不产生任何信号。因此，磁探针本质上是一个\textbf{磁场变化率探测器}。要获得静态磁场，需要主动积分（模拟RC积分电路或数字积分）。在实际托卡马克实验中，放电过程中磁场持续变化，探针可以直接测量 $\dot{B}$，再结合MHD平衡方程（第5章 Grad-Shafranov 方程）反推等离子体电流分布和压强分布。
\end{insight}

\subsection{Rogowski线圈：测量等离子体电流}

\begin{definition}{Rogowski线圈}{}{}
\textbf{Rogowski线圈}是将导线均匀绕制在截面恒定的非磁性环形骨架（通常为矩形或圆形截面）上，令被测电流穿过环中心。根据安培环路定律，线圈感应电压与穿过环的总电流变化率成正比：
\begin{equation}
\mathcal{E} = -\frac{\mu_0 N A}{2\pi R}\frac{\diff I}{\diff t},
\label{eq:rogowski}
\end{equation}
$R$ 为环的平均半径，$N$ 为总匝数，$A$ 为单匝截面积。对电压积分即可得到总电流 $I(t)$。
\end{definition}

\begin{insight}{Rogowski线圈的"自屏蔽"特性}{}{}
Rogowski线圈的一个重要优点是：它只测量\emph{穿过环中心的净电流}，对环外部电流不敏感。这是因为外部电流产生的磁场在环内部绕行一周时，从不同方向穿过绕组，贡献相互抵消。这一特性使得Rogowski线圈可以方便地套在真空室或导体棒外部，测量其中的等离子体电流，而不受外部杂散磁场的显著干扰。在托卡马克中，环向等离子体电流 $I_p$ 通常达到 $10^6\,\mathrm{A}$ 量级，Rogowski线圈缠绕在真空室外围，是监测这一核心MHD参数的标准手段。通过与极向磁探针阵列结合，还可以重构等离子体边界和内部的磁通面结构（安全因子 $q$ 分布），为等离子体平衡和稳定性分析提供输入。
\end{insight}

\begin{warning}{磁探针是否扰动等离子体？}{}{}
\textbf{误区：}磁探针只是被动接收磁场，完全不扰动等离子体。

\textbf{正解：}磁探针本身是导体，插入等离子体后会与等离子体发生相互作用。首先，探针表面可能形成鞘层（类似朗缪尔探针），改变局部电位；其次，探针作为固态障碍物，会扰动等离子体流，产生尾流或激波；第三，在快速变化的磁场中，探针中感应的涡电流可能产生次级磁场，干扰原始信号。因此，磁探针通常设计得尽可能小（直径 $1\sim 3\,\mathrm{mm}$），并使用高阻抗前置放大器以减少电流回路。Rogowski线圈则由于位于等离子体外部，扰动问题较轻，但需要注意杂散电容和分布电感对快速瞬态信号的影响。
\end{warning}


\section*{随堂自测}
\addcontentsline{toc}{section}{随堂自测}

\begin{miniquiz}
\begin{enumerate}
    \item 为什么朗缪尔探针的悬浮电位 $V_f$ 通常远低于等离子体电位 $V_p$？请从电子与离子的通量差异解释。
    
    \item 微波干涉诊断中，如果选择过高频率（如 $f = 300\,\mathrm{GHz}$）测量低密度辉光放电（$n_e \sim 10^{16}\,\mathrm{m}^{-3}$），会面临什么主要困难？
    
    \item 汤姆逊散射中，离子散射截面为什么可以忽略？如果离子质量突然变为 $m_i = 100\,m_e$（如某种假想粒子），散射信号会发生什么变化？
    
    \item 某光谱仪测得等离子体中两条谱线的强度比为 $I_1/I_2 = 5$，对应的能级差为 $\Delta E = 2\,\mathrm{eV}$，统计权重比 $g_1/g_2 = 2$。估算激发温度 $T_{\mathrm{exc}}$。（提示：利用式\eqref{eq:line-ratio}）
    
    \item 磁探针直接测量的是 $B$、$\dot{B}$ 还是 $\ddot{B}$？如果要测量静态磁场，需要增加什么处理步骤？
\end{enumerate}
\end{miniquiz}

\begin{review}{本节要点}{}{}
\begin{itemize}
    \item 等离子体诊断的本质是将理论预言（德拜长度、等离子体频率、分布函数等）转化为可测量信号。
    \item 朗缪尔探针通过 $I$-$V$ 曲线测量局部 $T_e$ 和 $n_e$；过渡区 $\ln I$-$V$ 斜率给出 $T_e$，饱和电流给出 $n_e$。
    \item 微波干涉通过相位移动 $\Delta\phi \propto n_e L$ 测量线平均密度；截止密度 $n_{\mathrm{cut}} \propto f^2$ 限制了可测密度上限。
    \item 汤姆逊散射光谱宽度直接反映电子热速度分布，是测量 $T_e$ 的"黄金标准"；离子散射可忽略因为 $(m_e/m_i)^2$ 极小。
    \item 光谱诊断通过线光谱（激发温度、成分）、连续谱（韧致辐射温度）和斯塔克展宽（密度）提供丰富的参数信息。
    \item 磁探针和Rogowski线圈基于法拉第感应，测量磁场变化和等离子体电流，是MHD宏观参数的主要诊断手段。
\end{itemize}
\end{review}

\begin{chapterreview}[第12章要点回顾：诊断——理论落地的桥梁]
\begin{enumerate}
    \item \textbf{诊断的分类与选择}：按原理分为电探针、电磁波、光学和粒子/磁探针四大类；按空间特性分为点测量（朗缪尔探针、汤姆逊散射）和线/面积分测量（微波干涉、光谱线积分）；按时间特性分为稳态和瞬态诊断。
    
    \item \textbf{朗缪尔探针}：单探针、双探针和发射探针的结构差异；$I$-$V$ 曲线的四个区域（离子饱和、过渡、电子饱和、悬浮）；电子温度 $T_e = [\diff(\ln I)/\diff V]^{-1}$；电子密度 $n_e \propto I_{\mathrm{e,sat}}/\sqrt{T_e}$；悬浮电位为负的本质原因；Druyvesteyn公式通过二阶导数求EEDF；磁化等离子体中需修正鞘层理论。
    
    \item \textbf{微波干涉}：相位移动 $\Delta\phi = (e^2/2\eps_0 m_e c\omega) \int n_e\,\diff l$；截止密度 $n_{\mathrm{cut}} = \eps_0 m_e \omega^2/e^2$；适用于 $n_e \gtrsim 10^{18}\,\mathrm{m}^{-3}$ 的高密度测量；频率选择需兼顾穿透能力与相位灵敏度。
    
    \item \textbf{激光汤姆逊散射}：非相干散射条件 $n_e \ll 10^{23}\,\mathrm{m}^{-3}$；散射光谱高斯宽度 $\Delta\omega/\omega_i \propto \sqrt{T_e}$；直接测量电子速度分布；离子散射截面可忽略；是聚变装置芯部 $T_e$ 测量的标准手段。
    
    \item \textbf{光谱诊断}：玻尔兹曼图法求激发温度；韧致辐射斜率求 $T_e$；斯塔克展宽 $\Delta\lambda \propto n_e^{2/3}$ 求密度；$H_\alpha$ 线常用于边界密度监测；需要精确原子物理数据支持定量反演。
    
    \item \textbf{磁探针与Rogowski线圈}：基于法拉第电磁感应 $\mathcal{E} = -N A\,\dot{B}$；Rogowski线圈测量总电流，具有外部电流自屏蔽特性；积分电路或数字积分用于从 $\dot{B}$ 恢复 $B$；是MHD平衡和稳定性分析的数据来源。
    
    \item \textbf{贯穿全章的核心思想}：诊断是理论与实验的桥梁。德拜长度（第2章）决定探针鞘层尺度；截止频率（第6章）决定微波传播特性；分布函数（第7章）决定散射光谱形状；碰撞与辐射过程（第8章）决定光谱强度；MHD方程（第5章）将磁探针数据与等离子体平衡联系起来。没有诊断，等离子体物理学将只是数学模型；没有理论，诊断数据只是无意义的电压和光强。
\end{enumerate}
\end{chapterreview}

\newpage
